\documentclass[UTF8,a4paper,11pt,fontset=fandol]{ctexart}
\usepackage[margin=22mm]{geometry}
\usepackage{amsmath,amssymb,booktabs,graphicx,float,xcolor,fancyhdr,fvextra,hyperref}
\hypersetup{colorlinks=true,linkcolor=blue!45!black,urlcolor=blue!45!black,pdftitle={SVD及其应用实验报告}}
\setlength{\parindent}{2em}
\setlength{\parskip}{4pt}
\linespread{1.18}
\pagestyle{fancy}\fancyhf{}\fancyhead[L]{科学计算实验}\fancyhead[R]{SVD 及其应用}\fancyfoot[C]{\thepage}
\setlength{\headheight}{15pt}
\DefineVerbatimEnvironment{code}{Verbatim}{fontsize=\small,breaklines=true,frame=single}
\begin{document}
\hypersetup{pageanchor=false}\begin{titlepage}
\centering\vspace*{22mm}
{\Large 科学计算课程\par}\vspace{12mm}
{\Huge\bfseries 第一节上机实验\par}\vspace{7mm}
{\LARGE SVD 及其应用\par}\vspace{18mm}
{\large 实验报告\par}\vspace{16mm}
\vfill
\end{titlepage}\hypersetup{pageanchor=true}
\tableofcontents\newpage
\section{实验目标、数据与完成范围}
按课程附件完成以下内容：给定矩阵的 SVD 与能量分析；灰度图重建、RGB 对比网格、全部 $k$ 的指标计算、三类曲线以及三个关键发现；共享基在人脸识别、低带宽上传、三维表示和时序编码中的提高实验。

图像采用 scikit-image 的 astronaut.png（NASA 图源），大小为 $512\times512\times3$；读入为 float64，并除以 255 归一化。人脸使用 AT\&T（ORL）公开数据，共 40 人、400 张 $112\times92$ 灰度图。原始图像与人脸数据保存在实验目录，重新运行无需下载。

\textbf{完成边界：}核心部分全部实现。Eigenfaces 与量化使用真实公开静态人脸；三维部分使用合成曲面，时序部分使用模拟轨迹，均属于教学机制验证，不代表真实三维扫描、真实视频或边缘设备部署结果。未虚构课堂分组讨论记录。
\section{给定矩阵的 SVD}
\subsection{矩阵组装与输入核对}
将附件矩阵逐行录入后发现，印出的 $A$ 对应 $Q_5\Sigma Q_4$，因此此处 $Q_4$ 扮演 $V^T$ 的角色。采用该乘法与印刷矩阵的最大元素差为 $3.276\times10^{-5}$；若使用 $Q_4^T$，差异达 41.700658。前者差异与输入仅保留六位小数相符。本报告统一分析组装矩阵，不擅自对输入重新正交化。

附件将 $Q$ 称为“正定矩阵”，但它们非对称且近似正交。实测 $\|Q_5^TQ_5-I\|_F=2.284\times10^{-6}$，$\|Q_4^TQ_4-I\|_F=2.176\times10^{-6}$。


\[A=\begin{pmatrix}-13.026561 & 0.285437 & 5.211156 & 3.325249\\-57.568127 & -14.043320 & 37.227236 & 46.963095\\21.533442 & 1.506950 & -10.469939 & -9.796990\\-28.005110 & -7.765161 & 18.982922 & 24.852919\\-8.844262 & -4.251763 & 7.694655 & 11.817520\end{pmatrix}.\]

\subsection{数学原理与算法}
对 $A\in\mathbb R^{m\times n}$，约化 SVD 写为
\[
 A=U\Sigma V^T=\sum_{i=1}^{r}\sigma_i u_i v_i^T,\qquad
 A_k=\sum_{i=1}^{k}\sigma_i u_i v_i^T.
\]
其中 $\sigma_i$ 降序排列，$u_i,v_i$ 为左右奇异向量，并满足
\[
 Av_i=\sigma_i u_i,\quad A^Tu_i=\sigma_i v_i,\quad
 A^TAv_i=\sigma_i^2v_i,\quad AA^Tu_i=\sigma_i^2u_i.
\]
非方阵 $A$ 不宜直接称有左右“特征向量”；这里的特征关系属于 $A^TA$ 和 $AA^T$。截断近似在 Frobenius 范数意义下最优，且
\[
 E(k)=\frac{\sum_{i=1}^k\sigma_i^2}{\sum_{i=1}^r\sigma_i^2},\qquad
 \frac{\|A-A_k\|_F}{\|A\|_F}=\sqrt{1-E(k)}.
\]
教学实现使用 \texttt{eigh} 求 $A^TA$ 特征对，再计算 $u_i=Av_i/\sigma_i$；以 \texttt{numpy.linalg.svd} 交叉核验。构造 $A^TA$ 会使条件数平方，因此后续计算直接使用 SVD。两种方法的奇异值最大绝对差为 $5.487\times10^{-11}$。
\subsection{秩、奇异值及能量结果}


\begin{center}\small\begin{tabular}{rrrr}\toprule{}$k$ & $\sigma_k$ & $E(k)$ & 相对误差\\\midrule
1 & 100.000015319 & 0.990098022738 & 9.951e-02\\
2 & 10.000004986 & 0.999999009805 & 9.951e-04\\
3 & 0.099999934 & 0.999999999901 & 9.950e-06\\
4 & 0.000999999 & 1.000000000000 & 5.811e-16\\\bottomrule\end{tabular}\end{center}

四个奇异值均为正，故 $\operatorname{rank}(A)=4$，为满列秩；$\kappa_2(A)\approx100000.071$。保留前两项已有约 $99.999901\%$ 能量，但精确秩仍为 4。$AA^T$ 有一个额外零特征值，$A^TA$ 的四个特征值均为正。附录给出左右奇异向量及四个秩一矩阵，求和重构误差约为 $5.975\times10^{-14}$。
\section{SVD 图像压缩}
\subsection{灰度图的三行核心代码}
\par\noindent\begin{minipage}{\linewidth}\begin{code}
A = img @ np.array([0.2126, 0.7152, 0.0722])
U, s, Vt = np.linalg.svd(A, full_matrices=False)
A_k = (U[:, :20] * s[:20]) @ Vt[:20, :]
\end{code}
\end{minipage}\par
输入 \texttt{img} 已转换为 $[0,1]$ 浮点数，输出形状为 $(512,512)$。显示时执行 \texttt{np.clip(A\_k,0,1)}。\par\texttt{full\_matrices=False} 返回约化形状，不等于只计算前 $k$ 项。


\begin{figure}[H]\centering\includegraphics[width=.90\linewidth]{figures/01_grayscale_rank20.png}\caption{灰度原图与 $k=20$ 的截断重建。}\end{figure}

\subsection{RGB 三通道与可视化网格}
分别对 R、G、B 分解，以相同 $k$ 截断后堆叠为彩色图。按附件选择
\[k\in\{1,2,3,5,8,12,20,35,50,80,120\}.\]
原图与 11 个重建组成 $3\times4$ 网格。图中可见小 $k$ 仅保留大致颜色和轮廓，随着 $k$ 增加，边缘与纹理逐步恢复。


\begin{figure}[H]\centering\includegraphics[width=1.0\linewidth]{figures/02_rgb_rank_grid.png}\caption{原图与各截断阶数的 RGB 重建；PSNR 使用裁剪前的理论误差。}\end{figure}

\subsection{不显式重建的指标计算}
设图像高、宽、通道数为 $H,W,C$，定义
\[
 T=\sum_c\sum_i\sigma_{c,i}^2,\qquad R_k=\sum_c\sum_{i>k}\sigma_{c,i}^2.
\]
由酉不变性，所有通道误差平方和就是 $R_k$，故
\[
 E(k)=1-\frac{R_k}{T},\quad e_{\mathrm{rel}}=\sqrt{\frac{R_k}{T}},\quad
 \mathrm{MSE}=\frac{R_k}{HWC},\quad\mathrm{PSNR}=10\log_{10}\frac{1}{\mathrm{MSE}}.
\]
\[
 P(k)=Ck(H+W+1),\qquad \mathrm{CR}(k)=\frac{HWC}{P(k)}=\frac{HW}{k(H+W+1)}.
\]
每通道分别保存 $Hk$ 个左向量元素、$k$ 个奇异值和 $kW$ 个右向量元素。满秩时尾和为 0，PSNR 记为 $+\infty$。从末尾累加奇异值平方，可避免总量减前缀导致的消减误差。
\par\noindent\begin{minipage}{\linewidth}\begin{code}
e = np.stack([s**2 for _, s, _ in svd_ch])
tail = np.concatenate([np.cumsum(e[:, ::-1], axis=1)[:, ::-1],
                       np.zeros((C, 1))], axis=1)
def stats(k):
    err2 = tail[:, k].sum()
    mse = err2 / (H * W * C)
    return (H * W / (k * (H + W + 1)),
            np.sqrt(err2 / e.sum()),
            np.inf if mse == 0 else 10 * np.log10(1 / mse))
\end{code}
\end{minipage}\par
实际程序另记录参数量、MSE 与能量。完整指标表包含 $k=1,\ldots,512$ 的全部 512 行，保存为 \path{image_metrics_all_k.csv}。下表为代表性阶数：


\begin{center}\small\begin{tabular}{rrrrr}\toprule{}$k$ & 能量（\%） & 相对误差 & PSNR（dB） & 参数压缩比\\\midrule
1 & 78.50571 & 0.463619 & 11.856 & 255.750\\
4 & 90.55658 & 0.307301 & 15.428 & 63.938\\
20 & 97.94973 & 0.143188 & 22.061 & 12.788\\
35 & 98.97123 & 0.101428 & 25.056 & 7.307\\
80 & 99.72259 & 0.052670 & 30.748 & 3.197\\
120 & 99.88258 & 0.034267 & 34.482 & 2.131\\
255 & 99.98974 & 0.010131 & 45.066 & 1.003\\
256 & 99.98992 & 0.010041 & 45.144 & 0.999\\
512 & 100.00000 & 0.000000 & $+\infty$ & 0.500\\\bottomrule\end{tabular}\end{center}

在 $k=1,4,20,80,120,256,512$ 上，程序将尾和结果与显式重建 MSE 交叉核验，全部通过。图像显示裁剪后的误差另存；裁剪改变矩阵，不能直接套用未裁剪的尾和等式。
\subsection{奇异值、质量与压缩比曲线}


\begin{figure}[H]\centering\includegraphics[width=.79\linewidth]{figures/03_singular_values.png}\caption{三通道奇异值衰减，纵轴为对数尺度。}\end{figure}

\begin{figure}[H]\centering\includegraphics[width=.79\linewidth]{figures/04_error_psnr.png}\caption{相对误差与 PSNR 随 $k$ 的变化；无穷值不绘作有限点。}\end{figure}

\begin{figure}[H]\centering\includegraphics[width=.79\linewidth]{figures/05_compression_ratio.png}\caption{参数压缩比及盈亏平衡点 $k^*=255.750$。}\end{figure}

\section{三个关键发现}
\subsection{能量占比不能替代视觉质量}
$k=1$ 时能量占比为 $78.5057\%$，PSNR 仅 11.856 dB；$k=4$ 时能量已有 $90.5566\%$，PSNR 仍为 15.428 dB。即使 $k=20$ 能量达到 $97.9497\%$，PSNR 也只有 22.061 dB。大面积背景、平均亮度贡献较大，能量较高不能保证纹理细节完整。
\subsection{盈亏平衡点：区分参数与字节}
由 $\mathrm{CR}>1$ 得
\[
 k<k^*=\frac{HW}{H+W+1}=\frac{262144}{1025}=255.750244.
\]
本图 $k\le255$ 时参数更少，$k\ge256$ 时参数更多。因此超过 $k^*$ 后，增加阶数虽然提高重建质量，但会“越压越大”。附件中约 353 对应 $620\times820$ 图像，不能沿用到本图。

以上假设每个标量占相同字节。原始 uint8 像素占 786432 B，$k=20$ 的 float32 因子占 246000 B，实际原始负载压缩比仅 $3.197$ 倍，而参数比为 $12.788$ 倍。float32 相对 uint8 的字节盈亏平衡点是 $255.750244/4\approx63.938$，从 $k=64$ 起因子已更大。PNG 文件大小为 791555 B，还包含编码与文件开销；本实验未与 JPEG 作率失真比较。
\subsection{细节失真观察}


\begin{figure}[H]\centering\includegraphics[width=1.0\linewidth]{figures/06_detail_crop.png}\caption{相同局部区域对照：原图、$k=1$、$k=4$、$k=20$。}\end{figure}

宇航服徽章、接缝和器件边缘较早模糊，大块颜色区域更容易保留。这表明本图细节需要更多奇异分量共同表达。小奇异值不普遍等同于高空间频率，观察到的趋势不是对所有图像成立的严格定理。
\section{提高实验：共享基表示}
\subsection{Eigenfaces 人脸识别}
将每张图展开为 $D=10304$ 维向量。固定随机种子 20260917，每人随机选 7 张训练、3 张测试，共 280/120 张。均值和基只用训练集计算，测试照片不参与拟合。两组包含相同身份、不同照片，属于已知身份的闭集识别。

令训练矩阵中心化为 $X_c$，对其 SVD 得到行基 $B_k$。投影与重建为
\[
 z=(x-\mu)B_k^T,\qquad \hat x=\mu+zB_k.
\]
识别在系数空间采用欧氏距离最近邻。主评估预先固定 $k=80$，其余 $k$ 用于敏感性分析，不以测试结果选择新参数。训练中心化矩阵数值秩为 279。


\begin{figure}[H]\centering\includegraphics[width=.90\linewidth]{figures/08_mean_eigenfaces.png}\caption{均值脸及前 11 个特征脸，各基图独立归一化显示。}\end{figure}

% Eigenfaces metrics: see CSV for all columns.

\begin{center}\small\begin{tabular}{rrrr}\toprule{}$k$ & 正确 / 120 & 准确率（\%） & 测试 PSNR（dB）\\\midrule
1 & 9 & 7.50 & 17.038\\
5 & 99 & 82.50 & 18.891\\
20 & 114 & 95.00 & 20.800\\
32 & 117 & 97.50 & 21.442\\
80 & 118 & 98.33 & 22.618\\
200 & 116 & 96.67 & 23.529\\
279 & 116 & 96.67 & 23.816\\\bottomrule\end{tabular}\end{center}

主结果为 118/120 正确，准确率 $98.33\%$；$k=200$ 和 279 时降为 $96.67\%$，虽然重建误差继续下降。这说明重建保真与识别判别力不是同一目标。单次划分结果不能解释为跨身份泛化保证。


\begin{figure}[H]\centering\includegraphics[width=.85\linewidth]{figures/07_eigenfaces_accuracy.png}\caption{识别准确率随保留主成分数量的变化。}\end{figure}

\begin{figure}[H]\centering\includegraphics[width=.90\linewidth]{figures/09_face_reconstructions.png}\caption{保留测试照片在不同阶数下的重建。}\end{figure}

共享基与均值也占存储。$N$ 张图的参数压缩比为
\[
 \frac{ND}{D+kD+Nk}.
\]
本训练库 $N=280,k=80$ 时为 $3.366$ 倍。即使保留训练集全部非零方向，测试图像仍可能具有子空间外分量，故测试重建不一定无误差。
\subsection{边缘上传：每维 int8 量化}
采用 $k=32$ 的训练基，以训练系数每维 $0.1\%$ 和 $99.9\%$ 分位数定义 $\mathrm{lo},\mathrm{hi}$，并令 $s=(\mathrm{hi}-\mathrm{lo})/255$。编码和解码为
\[
 q=\operatorname{round}\!\left(\operatorname{clip}\!\left(\frac{z-\mathrm{lo}}{s},0,255\right)\right)-128,
 \qquad \hat z=(q+128)s+\mathrm{lo}.
\]
量化器仅由训练集标定。均值、共享基、lo、scale 均作为模型保存；两端应校验基 ID、版本和量化参数。本实验按固定除以 255 归一化，不需要每帧传输额外亮度参数。


\begin{center}\small\begin{tabular}{ll}\toprule{}指标 & 本次结果\\\midrule
浮点 / int8 重建 PSNR & 21.442089 / 21.441348 dB\\
PSNR 损失 / 最大像素差 & 0.000742 dB / 0.021111\\
超范围系数比例 & $0.2344\%$\\
量化后识别准确率 & $96.67\%$\\
uint8 原帧 / 系数负载 & 10304 B / 32 B\\
共享模型原始负载 & 1360384 B\\
120 帧含模型压缩比 & 0.9064\\
摊销模型的盈亏平衡 & 133 帧\\\bottomrule\end{tabular}\end{center}

每帧负载虽缩小 322 倍，但短序列并不能抵消模型分发开销；相对原始像素节省需满足\[1360384+32N<10304N.\]以上未计网络包头、加密与重传。

平均 PSNR 损失很小，最大像素差却为 0.021111，不能复述附件中其他案例的 0.0015 为本次结果。系数化也不等于加密：系数仍能用于身份识别，持有基即可近似重建；未证明去标识化或隐私保护效果。
\subsection{三维共享基：合成曲面的教学演示}
3DMM 将对应顶点坐标和纹理分别表示为均值与基的线性组合：
\[
 S=\bar S+\sum_i\alpha_i b_i,\qquad T=\bar T+\sum_i\beta_i c_i.
\]
本次未获取真实三维扫描；构造 $40\times32$ 网格、1280 个顶点、3840 维坐标，以六种线性形变生成 200 个样本，160 个训练、40 个测试。只演示形状 PCA，不包含纹理 PCA 或从照片反演三维。


\begin{figure}[H]\centering\includegraphics[width=1.0\linewidth]{figures/11_synthetic_3dmm.png}\caption{合成曲面与低秩重建，非真实三维人脸扫描。}\end{figure}

\begin{center}\small\begin{tabular}{rr}\toprule{}$k$ & 测试曲面 RMSE\\\midrule
1 & 3.654625e-02\\
2 & 3.130014e-02\\
3 & 2.289169e-02\\
5 & 5.568749e-03\\
6 & 5.362108e-16\\\bottomrule\end{tabular}\end{center}

$k=6$ 时达到浮点精度，是因为数据生成机制本来只有六个线性自由度，不能据此声称真实人脸只需六个参数。
\subsection{时序共享基：模拟轨迹与差分编码}
ORL 没有连续采集时间戳，不能将文件顺序当视频。以一个测试样本的系数为中心，叠加振幅为训练系数标准差 0.12 倍的平滑正弦变化，生成 300 帧、32 维模拟轨迹。该构造主动满足平滑性，结果仅限于此条件。


\begin{figure}[H]\centering\includegraphics[width=.85\linewidth]{figures/10_synthetic_temporal_coefficients.png}\caption{模拟系数轨迹的前四维，非真实视频。}\end{figure}

每 60 帧记录绝对系数，其余保存相邻差；差分可能超出 int8，使用 int16，再用 zlib level=9 压缩。解压和累加后逐元素恢复量化系数，这不意味着对原图无损。


\begin{center}\small\begin{tabular}{lr}\toprule{}表示方式 & 字节数\\\midrule
直接 int8 系数 & 9600\\
int16 差分 & 19200\\
直接 int8 + zlib & 1979\\
int16 差分 + zlib & 1387\\\bottomrule\end{tabular}\end{center}

公平比较双方均使用 zlib，额外压缩为 $1979/1387=1.427$ 倍，未达到附件案例的 $2\sim3$ 倍。仅做差分而不熵编码，反而使负载翻倍。包含一次模型分发时，300 帧相对原始像素压缩比为 2.270 倍。

此处周期绝对系数帧仅重置预测状态，没有重算基。真实应用可监测重建残差或分布漂移来触发新基，并同步基版本；本实验未验证真实场景切换、漂移阈值或与视频编码器的竞争性能。
\section{结论与复现说明}
给定矩阵满列秩为 4，奇异值跨五个数量级，少量分量能保留主要能量。图像实验表明，增加阶数能提升重建质量，但同时增加存储；能量、视觉质量、参数量和文件字节必须区分。

共享基使图像压缩与识别使用同一低维表示。模型摊销、量化截断和时序编码都有实际代价，必须结合样本数量与数据类型评价。真实 ORL 实验与三维、时序教学模拟分别报告，不将模拟收益扩展为真实部署结论。

本次使用智能体辅助阅读材料、核对矩阵、整理程序、执行计算及生成报告。Python 程序与原始数据随提交目录附上，运行命令为：
\par\noindent\begin{minipage}{\linewidth}\begin{code}
python -m pip install -r 02_实验本体/requirements.txt
python 02_实验本体/run_experiment.py
python 03_LaTeX源码/build_report.py
bash 03_LaTeX源码/compile.sh
\end{code}
\end{minipage}\par


\begin{center}\small\begin{tabular}{ll}\toprule{}项目 & 本次运行环境\\\midrule
Python & 3.13.7\\
NumPy / SciPy & 2.5.3 / 1.18.1\\
Matplotlib / Pillow & 3.11.2 / 12.3.0\\
随机种子 & 20260917\\
自动检查 & 22/22 通过\\\bottomrule\end{tabular}\end{center}

完整指标、训练测试划分、逐张预测和校验项分别见实验结果目录中的 CSV 与 JSON；运行日志为 \path{run.log}。本报告由 LaTeX 编译，编译要求与复现步骤见文件夹说明。姓名、学号与班级未提供，封面留空。
\section*{参考资料}\addcontentsline{toc}{section}{参考资料}
\begin{enumerate}
\item 课程附件：《第一节上机实验——SVD 及其应用》，本次提交包保留原件。
\item AT\&T Laboratories Cambridge, The Database of Faces.\\\url{https://www.cl.cam.ac.uk/research/dtg/attarchive/facedatabase.html}
\item scikit-image, astronaut image.\\\url{https://scikit-image.org/docs/stable/api/skimage.data.html}
\item NumPy, numpy.linalg.svd.\\\url{https://numpy.org/doc/stable/reference/generated/numpy.linalg.svd.html}
\item Turk M, Pentland A. Eigenfaces for Recognition. Journal of Cognitive Neuroscience, 1991, 3(1):71--86. DOI: 10.1162/jocn.1991.3.1.71.
\item Blanz V, Vetter T. A Morphable Model for the Synthesis of 3D Faces. SIGGRAPH, 1999:187--194. DOI: 10.1145/311535.311556.
\end{enumerate}
\newpage\appendix
\section{左右奇异向量与秩一展开}
以下数值保留六位小数，完整精度数组见 \path{matrix_decomposition.npz}。$U$ 的第 $i$ 列为左奇异向量；$V^T$ 的第 $i$ 行转置为右奇异向量。左右向量成对变号不改变分解。


\[U=\begin{pmatrix}-0.130739 & 0.608762 & -0.043819 & -0.780233\\-0.842723 & -0.093659 & -0.312900 & 0.107165\\0.251167 & -0.638035 & -0.511983 & -0.505413\\-0.426363 & -0.217889 & 0.017690 & -0.144392\\-0.166897 & -0.407521 & 0.798582 & -0.321653\end{pmatrix},\qquad\]\[V^T=\begin{pmatrix}0.690420 & 0.161962 & -0.440611 & -0.550410\\-0.657117 & 0.395218 & -0.090602 & -0.635446\\0.008191 & 0.773622 & -0.356826 & 0.523563\\-0.302409 & -0.468060 & -0.818736 & 0.138346\end{pmatrix}.\]

令 $R_i=\sigma_i u_iv_i^T$，则 $A=R_1+R_2+R_3+R_4$，$A_k$ 为前 $k$ 项之和。

\[R_1=\begin{pmatrix}-9.026484 & -2.117475 & 5.760505 & 7.196006\\-58.183285 & -13.648908 & 37.131302 & 46.384313\\17.341075 & 4.067951 & -11.066695 & -13.824483\\-29.436956 & -6.905459 & 18.786022 & 23.467443\\-11.522902 & -2.703096 & 7.353665 & 9.186176\end{pmatrix}.\]

\[R_2=\begin{pmatrix}-4.000276 & 2.405937 & -0.551552 & -3.868355\\0.615446 & -0.370156 & 0.084857 & 0.595150\\4.192633 & -2.521629 & 0.578073 & 4.054368\\1.431787 & -0.861138 & 0.197412 & 1.384569\\2.677889 & -1.610597 & 0.369223 & 2.589578\end{pmatrix}.\]

\[R_3=\begin{pmatrix}-0.000036 & -0.003390 & 0.001564 & -0.002294\\-0.000256 & -0.024207 & 0.011165 & -0.016382\\-0.000419 & -0.039608 & 0.018269 & -0.026806\\0.000014 & 0.001369 & -0.000631 & 0.000926\\0.000654 & 0.061780 & -0.028495 & 0.041811\end{pmatrix}.\]

\[R_4=\begin{pmatrix}0.000236 & 0.000365 & 0.000639 & -0.000108\\-0.000032 & -0.000050 & -0.000088 & 0.000015\\0.000153 & 0.000237 & 0.000414 & -0.000070\\0.000044 & 0.000068 & 0.000118 & -0.000020\\0.000097 & 0.000151 & 0.000263 & -0.000044\end{pmatrix}.\]

六位小数表格会引入舍入误差，不能用表格精度要求复现机器精度残差。直接对指导书印刷 $A$ 分解所得奇异值约为 $100.000000439$、$10.000000130$、$0.100000094$、$0.001000092$；其与组装矩阵有微小输入差异，两种版本均保存在结果数组中。
\end{document}